\documentclass[10pt, aspectratio=169]{beamer}
% Preámbulo modular para presentaciones Beamer (Modelación Estadística)
% Diseño y paleta institucional del Tecnológico de Monterrey

\documentclass[aspectratio=169,xcolor={dvipsnames,table}]{beamer}

\usepackage[utf8]{inputenc}
\usepackage[spanish,mexico]{babel}
\usepackage{amsmath,amssymb,amsfonts,mathtools}
\usepackage{graphicx}
\usepackage{listings}
\usepackage{booktabs}
\usepackage{tikz}
\usetikzlibrary{matrix,arrows,decorations.pathmorphing,shapes.geometric,positioning}

% Paleta Institucional Tecnológico de Monterrey
\definecolor{TecAzul}{HTML}{0039A6}       % Azul TEC: color primario institucional
\definecolor{TecAzulOscuro}{HTML}{1A2E51} % Azul marino oscuro
\definecolor{TecRojo}{HTML}{EC2661}       % Rojo de acento (paleta miTec)
\definecolor{TecGris}{HTML}{646464}       % Gris neutro para comentarios y subtítulos
\definecolor{TecFondoBloque}{HTML}{F4F6F9}% Fondo suave para bloques

% Configuración de Tema Beamer
\usetheme{Boadilla}
\usecolortheme[named=TecAzulOscuro]{structure}
\setbeamercolor{alerted text}{fg=TecRojo}
\setbeamercolor{palette primary}{bg=TecAzulOscuro,fg=white}
\setbeamercolor{palette secondary}{bg=TecAzul,fg=white}
\setbeamercolor{palette tertiary}{bg=TecAzulOscuro!80!black,fg=white}
\setbeamercolor{title}{fg=TecAzulOscuro,bg=TecFondoBloque}
\setbeamercolor{frametitle}{fg=TecAzulOscuro,bg=TecFondoBloque}

% Bloques personalizados elegantes
\setbeamercolor{block title}{bg=TecAzulOscuro,fg=white}
\setbeamercolor{block body}{bg=TecFondoBloque,fg=black}
\setbeamercolor{block title alerted}{bg=TecRojo,fg=white}
\setbeamercolor{block body alerted}{bg=TecRojo!8,fg=black}
\setbeamercolor{block title example}{bg=TecAzul,fg=white}
\setbeamercolor{block body example}{bg=TecAzul!8,fg=black}

% Redondear bordes y añadir sombra sutil
\setbeamertemplate{blocks}[rounded][shadow=true]
\setbeamertemplate{navigation symbols}{}

% Pie de página limpio con número de diapositiva
\setbeamertemplate{footline}{
  \leavevmode%
  \hbox{%
  \begin{beamercolorbox}[wd=.7\paperwidth,ht=2.25ex,dp=1ex,left]{author in head/foot}%
    \hspace*{2ex}\insertshortauthor~~(\insertshortinstitute) \hfill \insertshorttitle
  \end{beamercolorbox}%
  \begin{beamercolorbox}[wd=.3\paperwidth,ht=2.25ex,dp=1ex,right]{date in head/foot}%
    \insertshortdate{}\hspace*{2em}
    \insertframenumber{} / \inserttotalframenumber\hspace*{2ex}
  \end{beamercolorbox}}%
  \vskip0pt%
}

% Configuración de Listings de Python para visualización pedagógica de código
\lstset{
    language=Python,
    basicstyle=\ttfamily\footnotesize,
    keywordstyle=\color{TecAzulOscuro}\bfseries,
    stringstyle=\color{TecRojo},
    commentstyle=\color{TecGris}\itshape,
    showstringspaces=false,
    breaklines=true,
    frame=lines,
    rulecolor=\color{TecAzulOscuro!30},
    backgroundcolor=\color{TecFondoBloque},
    aboveskip=0.5em,
    belowskip=0.5em,
    literate={á}{{\'a}}1 {é}{{\'e}}1 {í}{{\'i}}1 {ó}{{\'o}}1 {ú}{{\'u}}1
             {Á}{{\'A}}1 {É}{{\'E}}1 {Í}{{\'I}}1 {Ó}{{\'O}}1 {Ú}{{\'U}}1
             {ñ}{{\~n}}1 {Ñ}{{\~N}}1 {¿}{{?`}}1 {¡}{{!`}}1
}

% Metadatos institucionales por defecto
\title[Modelación Estadística]{Modelación Estadística para Ciencia de Datos e Ingeniería}
\author[J. Castillo Colmenares]{Juliho David Castillo Colmenares}
\institute[Tec de Monterrey]{Tecnológico de Monterrey \\ Escuela de Ingeniería y Ciencias}
\date{\today}

% Comandos y macros estadísticas compartidas para las presentaciones Beamer (Metropolis)

% Conjuntos numéricos básicos
\newcommand{\R}{\mathbb{R}}
\newcommand{\N}{\mathbb{N}}
\newcommand{\Q}{\mathbb{Q}}
\newcommand{\Z}{\mathbb{Z}}
\newcommand{\set}[1]{\left\{ #1 \right\}}
\newcommand{\abs}[1]{\left|#1\right|}
\newcommand{\norm}[1]{\left\| #1 \right\|}

% Comandos estadísticos y probabilísticos del curso
\newcommand{\Var}{\operatorname{Var}}
\newcommand{\cov}{\operatorname{Cov}}
\newcommand{\corr}{\rho}
\newcommand{\comb}[2]{\begin{pmatrix} #1 \\ #2 \end{pmatrix}}
\newcommand{\s}{\sigma}
\newcommand{\particion}{\mathfrak{P}}
\newcommand{\card}[1]{\#\left( #1 \right)}
\newcommand{\seq}[3]{\set{#1}_{#2}^{#3}}
\newcommand{\E}{\mathbb{E}}
\newcommand{\Prob}{\mathbb{P}}

% Entornos pedagógicos y cajas en español académico natural
\newenvironment{discusion}{%
  \begin{alertblock}{Pregunta de Discusión}%
}{%
  \end{alertblock}%
}

\newenvironment{reflexion}{%
  \begin{alertblock}{Reflexión para el Aula}%
}{%
  \end{alertblock}%
}

\newenvironment{paradoja}{%
  \begin{alertblock}{Paradoja de Exploración}%
}{%
  \end{alertblock}%
}

\newenvironment{motivacion}{%
  \begin{exampleblock}{Motivación: Ciencia de Datos en la Práctica}%
}{%
  \end{exampleblock}%
}

\newenvironment{retoclase}{%
  \begin{block}{Reto Rápido en Equipo}%
}{%
  \end{block}%
}


\title[Introducción a la Regresión]{Sección 09.02: Introducción a la Regresión Lineal}
\subtitle{Modelos Estocásticos, Supuestos Clásicos, y Aplicaciones}
\author[J. Castillo Colmenares]{Juliho Castillo Colmenares}
\institute{Tecnológico de Monterrey}
\date{\vspace{-1.2cm}}

\begin{document}

% Slide 1: Portada
\begin{frame}[plain]
  \titlepage
\end{frame}

% Slide 2: Hoja de Ruta
\begin{frame}{Hoja de Ruta --- Capítulo 09: Regresiones Lineales y Múltiples}
  \scriptsize
  ¿Dónde nos encontramos en el desarrollo modular del capítulo?
  \begin{itemize}\itemsep=0.02em
    \item \textbf{09.01} Correlación como Premisa de la Regresión
    \item \textbf{09.02} Introducción a la Regresión Lineal \emph{(Hoy)}
    \item \textbf{09.03} Matemáticas de la Regresión: Derivación MCO y $R^2$
    \item \textbf{09.04} Regresión Lineal sobre Datos Simulados
    \item \textbf{09.05} Coeficientes Óptimos, Pruebas $t$/$F$ y RSE
    \item \textbf{09.06} Implementación en Python con \texttt{statsmodels}
    \item \textbf{09.07} Regresión Múltiple, Ecuación Normal y Ridge/Lasso
    \item \textbf{09.08} Validación de Modelos y $k$-fold Cross-Validation
    \item \textbf{09.09} Diagnóstico de Residuos y Multicolinealidad (VIF)
    \item \textbf{09.10} Regresión con \texttt{scikit-learn}
    \item \textbf{09.11} Variables Categóricas y Variables Muda
    \item \textbf{09.12} Transformaciones No Lineales y Regresión Polinomial
  \end{itemize}
  \pause
  \vspace{-0.05cm}
  \begin{block}{Objetivo de la Sesión}
    Presentar la regresión lineal como modelo estocástico predictivo: distinguirla de los modelos determinísticos, diferenciarla formalmente de la correlación (Sección 09.01), enumerar sus cinco supuestos clásicos, y ubicar la regresión simple y múltiple dentro del panorama de aplicaciones en ciencia de datos.
  \end{block}
\end{frame}

% Slide 3: Motivación
\begin{frame}{De la Asociación al Modelo Predictivo}
  \footnotesize
  \begin{motivacion}
    La regresión lineal es una técnica desarrollada en el siglo XIX por Francis Galton y formalizada matemáticamente por Karl Pearson y Ronald Fisher. A pesar de su simplicidad, sirve de base para técnicas más avanzadas: regresión polinomial, modelos lineales generalizados, y métodos de aprendizaje automático.
  \end{motivacion}

  \vspace{0.15cm}
  \pause
  \begin{itemize}\itemsep=0.1cm
    \item \textbf{Predicción:} estimar valores futuros a partir de datos históricos. \pause
    \item \textbf{Inferencia:} comprender y cuantificar relaciones entre variables.
  \end{itemize}
\end{frame}

% Slide 4: Modelos determinísticos vs. estocásticos
\begin{frame}{Modelos Determinísticos vs. Estocásticos}
  \footnotesize
  \begin{block}{Modelo Determinístico}
    Produce siempre el mismo resultado para las mismas entradas (ej. $F=ma$): no hay término de error.
  \end{block}
  \pause

  \vspace{0.1cm}
  \begin{alertblock}{Modelo Estocástico}
    La regresión lineal reconoce que las observaciones reales no siguen perfectamente una relación lineal: incluye un componente de error aleatorio $\epsilon$, reflejando factores no capturados por el modelo.
  \end{alertblock}
\end{frame}

% Slide 5: Ejemplo del modelo
\begin{frame}{Ejemplo: el Precio de una Casa como Modelo Lineal}
  \footnotesize
  \begin{block}{Especificación del Modelo}
    Si el precio $P$ depende linealmente del tamaño $S$, las comodidades $A$ y el acceso a transporte $T$:
    $$P = a_1 S + a_2 A + a_3 T + \epsilon$$
  \end{block}
  \pause

  \vspace{0.1cm}
  \begin{alertblock}{Componentes}
    $P$ es la variable de salida (predicha); $S,A,T$ son las variables de entrada (conocidas); $a_1,a_2,a_3$ son parámetros que deben \textbf{estimarse a partir de datos históricos}; $\epsilon$ es el error aleatorio irreducible.
  \end{alertblock}
\end{frame}

% Slide 6: Correlación vs. regresión
\begin{frame}{Correlación vs. Regresión: la Distinción Formal}
  \footnotesize
  \begin{block}{Correlación (Sección 09.01)}
    Mide fuerza y dirección; es \textbf{simétrica} ($r_{XY}=r_{YX}$) y adimensional.
  \end{block}
  \pause

  \vspace{0.1cm}
  \begin{alertblock}{Regresión}
    Modela la relación funcional para \textbf{predecir} valores; \textbf{no es simétrica} (regredir $Y$ sobre $X$ da coeficientes distintos que regredir $X$ sobre $Y$) y sus coeficientes tienen unidades.
  \end{alertblock}
\end{frame}

% Slide 7: Supuestos clásicos
\begin{frame}{Los Cinco Supuestos Clásicos (Vista Previa)}
  \footnotesize
  Para que las inferencias de un modelo de regresión sean válidas:
  \begin{enumerate}\itemsep=0.06cm
    \item \textbf{Linealidad} de la relación predictor-respuesta. \pause
    \item \textbf{Independencia} de las observaciones. \pause
    \item \textbf{Homocedasticidad}: varianza constante del error. \pause
    \item \textbf{Normalidad} de los errores. \pause
    \item \textbf{No multicolinealidad} entre predictores.
  \end{enumerate}
  \vspace{0.05cm}
  \begin{alertblock}{Nota}
    El diagnóstico formal completo de estos cinco supuestos se desarrolla en la Sección 09.09.
  \end{alertblock}
\end{frame}

% Slide 8: Tipos de regresión
\begin{frame}{Regresión Simple vs. Regresión Múltiple}
  \footnotesize
  \begin{block}{Regresión Lineal Simple}
    Una sola variable predictora: $Y=\beta_0+\beta_1X+\epsilon$ (Secciones 09.03-09.06).
  \end{block}
  \pause

  \vspace{0.1cm}
  \begin{alertblock}{Regresión Lineal Múltiple}
    Dos o más predictores: $Y=\beta_0+\beta_1X_1+\cdots+\beta_pX_p+\epsilon$ (Sección 09.07). Ambos tipos comparten el mismo principio: mínimos cuadrados ordinarios.
  \end{alertblock}
\end{frame}

% Slide 9: Lab Python (Bloque 1)
\begin{frame}[fragile]{Laboratorio en Python: Determinístico vs. Estocástico}
  \scriptsize
  Varianza residual nula en un modelo determinístico vs. varianza no nula en uno estocástico (\texttt{09.02\_introduction\_to\_regression.py}):
  {\lstset{basicstyle=\fontsize{3.6pt}{4.3pt}\selectfont\ttfamily,aboveskip=0.05em,belowskip=0.05em}
  \lstinputlisting[firstline=17, lastline=27]{../../code/09_regresiones/09.02_introduction_to_regression.py}}
\end{frame}

% Slide 10: Lab Python (Bloque 2)
\begin{frame}[fragile]{Laboratorio en Python: Regresión Simple vs. Múltiple}
  \scriptsize
  Comparación de $R^2$ al pasar de un predictor a tres, usando el ejemplo del precio de la vivienda:
  {\lstset{basicstyle=\fontsize{3.2pt}{3.9pt}\selectfont\ttfamily,aboveskip=0.05em,belowskip=0.05em}
  \lstinputlisting[firstline=30, lastline=47]{../../code/09_regresiones/09.02_introduction_to_regression.py}}
\end{frame}

% Slide 11: Lab Python (Bloque 3)
\begin{frame}[fragile]{Laboratorio en Python: Vista Previa de los Supuestos}
  \scriptsize
  Verificación numérica rápida de media-cero y homoscedasticidad (auditoría completa en 09.09):
  {\lstset{basicstyle=\fontsize{3.4pt}{4.1pt}\selectfont\ttfamily,aboveskip=0.05em,belowskip=0.05em}
  \lstinputlisting[firstline=50, lastline=64]{../../code/09_regresiones/09.02_introduction_to_regression.py}}
\end{frame}

% Slide 12: Salida Terminal
\begin{frame}[fragile]{Laboratorio en Python: Salida en Terminal}
  \scriptsize
  Salida estándar generada al ejecutar el script del laboratorio en Python:
  \vspace{0.02cm}
  \begin{block}{Salida Estándar en Consola}
    \fontsize{3.6pt}{4.3pt}\selectfont
    \begin{verbatim}
=== Block 1: Deterministic vs. Stochastic ===
Deterministic (F=ma):     residual variance = 0.000000
Stochastic (price~size):  residual variance = 65.9973

=== Block 2: Simple vs. Multiple Regression ===
Simple (price~size):                  R^2=0.6728
Multiple (price~size+amenities+transit): R^2=0.8893 (+0.2165)

=== Block 3: Assumptions Quick Preview ===
mean(residuals) = -0.000000 (zero-mean OK)
Var(first half)=2.35, Var(second half)=2.45 (homoscedasticity OK)
    \end{verbatim}
  \end{block}
\end{frame}

% Slide 13: Aplicaciones
\begin{frame}{Aplicaciones Prácticas de la Regresión Lineal}
  \footnotesize
  \begin{itemize}\itemsep=0.1cm
    \item \textbf{Finanzas:} predicción de precios de acciones, evaluación de riesgo crediticio. \pause
    \item \textbf{Marketing:} modelar el impacto de la publicidad en ventas (el caso de la Sección 09.01). \pause
    \item \textbf{Medicina:} relacionar dosis de medicamentos con respuestas clínicas. \pause
    \item \textbf{Ingeniería y ciencias sociales:} control de calidad, predicción de fallas, estudios socioeconómicos.
  \end{itemize}
\end{frame}

% Slide 14: Cierre
\begin{frame}{Cierre de Sección: Introducción a la Regresión Lineal}
  \begin{reflexion}
    La regresión lineal formaliza la asociación medida por la correlación en un modelo estocástico predictivo: reconoce explícitamente que las observaciones reales incluyen error aleatorio, y descansa sobre cinco supuestos clásicos que garantizan la validez de sus inferencias.
  \end{reflexion}

  \vspace{0.2cm}
  \pause
  \begin{block}{Lecciones Clave}
    \begin{itemize}\itemsep=0.08cm
      \item \textbf{Estocástico, no determinístico:} el término $\epsilon$ es central, no un detalle menor.
      \item \textbf{Correlación $\neq$ regresión:} la regresión predice y sus coeficientes tienen unidades e interpretación asimétrica.
      \item \textbf{Simple vs. múltiple:} mismo principio (mínimos cuadrados), distinto número de predictores.
    \end{itemize}
  \end{block}
\end{frame}

% Slide 15: Perspectiva Modular
\begin{frame}{Perspectiva Modular --- El Próximo Reto}
  \scriptsize
  \begin{retoclase}
    Con el modelo conceptualmente establecido, la Sección 09.03 responde la pregunta matemática central: ¿cómo se estiman exactamente los parámetros $\beta_0,\beta_1$ a partir de los datos? Derivaremos rigurosamente las fórmulas de mínimos cuadrados ordinarios (MCO) y demostraremos la descomposición fundamental $\text{SCT}=\text{SCR}+\text{SCE}$ que da origen al coeficiente $R^2$.
  \end{retoclase}

  \vspace{0.08cm}
  \pause
  \begin{alertblock}{Preparación Recomendada}
    \begin{itemize}\itemsep=0.06cm
      \item Repasa el cálculo diferencial básico: mínimos de una función mediante derivadas parciales.
      \item Reflexiona sobre qué significa "minimizar la suma de errores cuadráticos" geométricamente.
    \end{itemize}
  \end{alertblock}
\end{frame}

\end{document}
